\documentclass[aspectratio=169]{beamer}

\usetheme{Madrid}
\usecolortheme{seagull}
\setbeamertemplate{navigation symbols}{}
\setbeamertemplate{footline}[frame number]

\usepackage[utf8]{inputenc}
\usepackage[T1]{fontenc}
\usepackage{booktabs}
\usepackage{array}
\usepackage{hyperref}

\title[LLM Stock Agent]{LLM-Assisted Multi-Agent Stock Decision Support System}
\subtitle{Final Course Project Presentation}
\author{Mo Di (2021030019)}
\institute{Representation and Reasoning: Foundations of Large Models}
\date{Final Presentation}

\begin{document}

\begin{frame}
  \titlepage
\end{frame}

\begin{frame}{Problem and Goal}
  \begin{itemize}
    \item Build an applied prototype for daily-frequency simulated stock decision support.
    \item Support three tasks: single-stock analysis, candidate screening, and portfolio rebalancing.
    \item Avoid a black-box ``ask the LLM for a stock pick'' workflow.
    \item Produce reproducible outputs, explicit risk warnings, and benchmarkable behavior.
  \end{itemize}
\end{frame}

\begin{frame}{System Workflow}
  \begin{center}
  \begin{tabular}{>{\centering\arraybackslash}p{0.20\linewidth}
                  >{\centering\arraybackslash}p{0.20\linewidth}
                  >{\centering\arraybackslash}p{0.20\linewidth}
                  >{\centering\arraybackslash}p{0.20\linewidth}}
    \toprule
    Data Layer & Agent Layer & Decision Layer & Evaluation Layer \\
    \midrule
    Prices, news, fundamentals, constraints
    &
    Market, sentiment, fundamental, risk agents
    &
    Decision synthesis plus optional OpenRouter reviewer
    &
    Fixed cases, baselines, backtest, dashboard \\
    \bottomrule
  \end{tabular}
  \end{center}
  \vspace{0.6em}
  \begin{block}{Design principle}
    The LLM sees structured intermediate evidence and reviews the decision; deterministic fallback keeps the benchmark reproducible.
  \end{block}
\end{frame}

\begin{frame}{Agent Design}
  \begin{columns}
    \begin{column}{0.48\linewidth}
      \begin{block}{Signal agents}
        \begin{itemize}
          \item MarketInformationAgent
          \item NewsSentimentAgent
          \item FundamentalAgent
        \end{itemize}
      \end{block}
      \begin{block}{Control agent}
        \begin{itemize}
          \item RiskManagerAgent checks volatility, drawdown, concentration, cash, and trade limits.
        \end{itemize}
      \end{block}
    \end{column}
    \begin{column}{0.48\linewidth}
      \begin{block}{Decision agents}
        \begin{itemize}
          \item DecisionAgent combines weighted evidence.
          \item LLMDecisionAdvisor optionally adjusts the result through OpenRouter JSON review.
        \end{itemize}
      \end{block}
      \begin{block}{Output}
        Action, score, confidence, rationale, risk warnings, target weights, and trade deltas.
      \end{block}
    \end{column}
  \end{columns}
\end{frame}

\begin{frame}{Data and Benchmark}
  \begin{itemize}
    \item Offline processed dataset: 12 large-cap tickers, 60 trading days.
    \item Inputs: OHLCV price history, sample news events, fundamentals, and portfolio constraints.
    \item Final benchmark: 15 fixed cases.
  \end{itemize}
  \vspace{0.4em}
  \begin{center}
  \begin{tabular}{lrr}
    \toprule
    Task & Cases & Purpose \\
    \midrule
    Single-stock analysis & 10 & action, rationale, risk awareness \\
    Candidate screening & 2 & ranked stock pool selection \\
    Rebalancing & 3 & portfolio constraints and trade deltas \\
    \bottomrule
  \end{tabular}
  \end{center}
\end{frame}

\begin{frame}{Evaluation Results}
  \begin{center}
  \begin{tabular}{lr}
    \toprule
    Metric & Result \\
    \midrule
    Unit tests & 6/6 passing \\
    Final benchmark cases & 15/15 passing \\
    Multi-agent risk warnings & 16 \\
    Equal-weight cumulative return & 4.77\% \\
    Equal-weight max drawdown & -1.37\% \\
    Equal-weight Sharpe ratio & 2.44 \\
    \bottomrule
  \end{tabular}
  \end{center}
  \vspace{0.5em}
  \begin{block}{Interpretation}
    The benchmark evaluates structure, constraints, and reproducibility. It is not a claim of live trading profitability.
  \end{block}
\end{frame}

\begin{frame}[fragile]{Demo Plan}
  \begin{enumerate}
    \item Run the final benchmark:
{\scriptsize
\begin{verbatim}
.\.venv\Scripts\python.exe scripts\run_demo.py --task all --llm off
\end{verbatim}
}
    \item Open generated outputs:
      \begin{itemize}
        \item \texttt{reports/final\_benchmark\_results.md}
        \item \texttt{reports/final\_dashboard.html}
      \end{itemize}
    \item Optional LLM mode:
{\scriptsize
\begin{verbatim}
$env:OPENROUTER_API_KEY="..."
.\.venv\Scripts\python.exe scripts\run_demo.py --task final --llm required
\end{verbatim}
}
  \end{enumerate}
\end{frame}

\begin{frame}{Engineering Deliverables}
  \begin{itemize}
    \item Runnable Python package under \texttt{src/stock\_agent/}.
    \item CLI entry point: \texttt{scripts/run\_demo.py}.
    \item Deterministic data generator: \texttt{scripts/generate\_final\_data.py}.
    \item Automated tests under \texttt{tests/}.
    \item Reports: Markdown summary and static HTML dashboard.
    \item Final written report and this presentation deck.
  \end{itemize}
\end{frame}

\begin{frame}{Limitations and Next Steps}
  \begin{itemize}
    \item Current data is deterministic and offline; live market/news feeds are future work.
    \item Scoring rules are transparent but simple; calibrated statistical models could improve them.
    \item The backtest is auxiliary and equal-weight, not an optimized trading strategy.
    \item Next steps: real OpenRouter run record, richer demo video, broader stock universe, and more realistic transaction-cost modeling.
  \end{itemize}
\end{frame}

\begin{frame}{Takeaway}
  \begin{block}{Final outcome}
    The project has moved from a minimal prototype to a complete applied-track deliverable: runnable system, fixed benchmark, baselines, tests, generated reports, final report, and presentation slides.
  \end{block}
\end{frame}

\end{document}
